\section{\label{sec:wc}sOPE 的威尔逊系数}

计算涂抹威尔逊系数的手续与 OPE 的情况平行，后者可参见例如
Ref.~\cite{Collins:1984rdg}。有了第~\ref{sec:phi4}~节的费曼规则，
把涂抹威尔逊系数计算到次领头阶是直截了当的。我们从非连通算符出发，确定领头阶连通与非连通算符的
涂抹威尔逊系数。

\subsection{非连通贡献}

\subsubsection{领头阶}

我们在图~\ref{fig:sope2pt} 中给出单位算符 $\mathbb{I}$ 这一非连通算符的涂抹威尔逊系数的
领头阶（树图级）与次领头阶（单圈）贡献。
\begin{figure}[htbp]
\begin{minipage}{0.49\textwidth}
\includegraphics[width=0.95\textwidth,keepaspectratio=true]{images/d1tree}
\end{minipage}%
\begin{minipage}{0.49\textwidth}
\includegraphics[width=0.95\textwidth,keepaspectratio=true]{images/d1loop}
\end{minipage}
\caption{\label{fig:sope2pt}表示威尔逊系数 $d_{\mathbb{I}}$ 在领头阶（左两图）与次领头阶
（右两图）之贡献的费曼图。黑色实线与蓝色虚线分别是零流时间与非零流时间处的传播子。黑色方块是未
涂抹场 $\phi(0)$，黑点是零流时间处的相互作用顶点，蓝色圆斑表示涂抹算符 $\rho^2(\tau,0)$。}
\end{figure}
正如我们将要看到的，领头阶贡献与重正化尺度 $\mu$ 无关，因为在这一阶流时间充当紫外调节子。
然而超出领头阶，涂抹威尔逊系数将出现重正化尺度依赖。

我们通过考虑 sOPE 中每个算符在真空中的矩阵元来提取非连通的涂抹威尔逊系数；这会消去一切连通贡献。
到 ${\cal O}(x)$ 精度，我们有
\begin{align}
\big\langle \Omega  |  {} &{\cal T}\{
\phi(x) \phi(0)\}  | \Omega \big\rangle =
\frac{d_{\mathbb{I}}(\mu x,\mu^2\tau,mx)}{4\pi^2x^2}\big\langle \Omega | \,
\mathbb{I} \,| \Omega \big\rangle
\nonumber\\
{} & +
d_{\rho^2}(\mu x,\mu^2\tau,mx)\big\langle \Omega | [\rho^2(\tau,0)
]_{\mathrm{R}}| \Omega
\big\rangle  +
{\cal O}(x),
\end{align}
这里我们对 $d_{\mathbb{I}}$ 的归一化作了选取，使其在自由理论中领头阶为 1。我们把该表达式中的
每个量按
\begin{equation}\label{eq:drhoexp}
f = f^{(0)} - \lambda f^{(1)} + {\cal O}(\lambda^2),
\end{equation}
展开到单圈，其中 $f$ 代表某个矩阵元或某个威尔逊系数。我们这样选取单圈项的符号，
是为了提出四点顶点费曼规则 $V^{(4)} = -\lambda/24$ 中耦合常数的符号。

于是，树图级的威尔逊系数由如下表达式的小时空行为给出：
\begin{align}
\frac{d_{\mathbb{I}}^{(0)}}{4\pi^2x^2}  \stackrel{x\sim 0}{=}{} &
\Big\{\big\langle
\Omega |
{\cal T}\{
\phi(x){} \phi(0)\} | \Omega \big\rangle \nonumber\\
{} & \qquad - \big\langle \Omega |
\rho^2(\tau,0)  | \Omega
\big\rangle\Big\}_{{\cal O}(m^2)}^{(0)} ,
\end{align}
其中为清晰起见我们略去了系数的自变量，并利用了树图级连通系数为 $d_{\rho^2}^{(0)}=1$
这一事实。下标表明我们必须把结果展开到 ${\cal O}(m^2)$。关于威尔逊系数计算的更多细节可参见，
例如，Ref.~\cite{Collins:1984rdg}。
相应的费曼积分表示为
\begin{align}
d_{\mathbb{I}}^{(0)} \stackrel{x\sim 0}{=} {} & 4\pi^2 x^2\left\{\int_k
\frac{e^{ik
\cdot
x}-e^{-2k^2 \tau
}}{k^2+m^2}\right\}_{{\cal O}(m^2)}
\nonumber \\
= {} & 4\pi^2x^2\int_k \left(e^{ik \cdot x} - e^{-2k^2
\tau}\right)\left(\frac{1}{k^2} - \frac{m^2}{(k^2)^2}\right),
\end{align}
其中
\begin{equation}
\int_k \equiv \int\frac{\mathrm{d}^4k}{(2\pi)^4}.
\end{equation}

涂抹威尔逊系数为
\begin{equation}\label{eq:d1lead}
d_{\mathbb{I}}^{(0)} = 1 -
\frac{x^2}{8\tau}  +\frac{m^2x^2}{4}\left[\gamma_{\;\mathrm{E}} -1 +
\log\left(\frac{x^2}{8\tau}\right)\right],
\end{equation}
其中 $\gamma_{\;\mathrm{E}}\simeq 0.577216$ 是欧拉--马歇罗尼常数。

我们可以把这个结果与 $d = 4-2\epsilon$
维时 OPE 的威尔逊系数 \cite{Collins:1984rdg} 相比较：
\begin{equation}
\overline{c}_{\mathbb{I}}^{(0)} =
1 + \frac{
m^2x^2}{4} \left[\frac{1}{\epsilon} + 1 +\gamma_{\;\mathrm{E}} +
\log\left(\frac{\pi\mu^2x^2}{4}\right)\right].
\end{equation}
于是在 $\ms$ 方案中，它变为
\begin{equation}
\overline{c}_{\mathbb{I}}^{(0)} =
1 +\frac{
m^2x^2}{4} \left[1 +  2\gamma_{\;\mathrm{E}} +
\log\left(\frac{\mu^2x^2}{16}\right)\right].
\end{equation}

尽管这些表达式中的有限部分无法直接比较——因为我们是在两种不同的重正化方案下表述威尔逊系数——
我们注意到三个重要特征。其一，对时空间隔的对数依赖完全相同。其二，在领头阶，流时间 $\tau$
扮演着重正化尺度的角色。其三，我们看到对小时间空间隔需要小的流时间参数。如果不恰当地选择流时间
参数，就会在涂抹威尔逊系数中产生大的对数贡献，即使在小时间空间隔下 sOPE 也会表现出糟糕的收敛性质。

\subsubsection{次领头阶}

单圈时，涂抹威尔逊系数由
\begin{align}
-\lambda\frac{d_{\mathbb{I}}^{(1)}}{4\pi^2x^2} = {} & \Big\{\big\langle
\Omega |
{\cal T}\{
\phi(x){} \phi(0)\} | \Omega \big\rangle \nonumber\\
{} & \qquad \qquad  - \big\langle \Omega |
\rho^2(\tau,0)  | \Omega
\big\rangle\Big\}_{{\cal O}(m^2)}^{(1)}
\end{align}
给出。

此图中的四点相互作用（示于图~\ref{fig:sope2pt}）出现在零流时间处。因此流时间无法充当对
$k_2$ 动量积分的调节子，我们必须引入重正化手续。我们采用维数正则化与 $\ms$ 方案。不过这个双重积分
很直接，因为两个积分可以分开进行：
\begin{equation}\label{eq:d11loop}
\frac{d_{\mathbb{I}}^{(1)}}{4\pi^2x^2}=  \bigg\{\int_{k_1}
\frac{e^{ik_1 \cdot
x} - e^{- 2k_1^2\tau}}{(k_1^2+m^2)^2}
\frac{1}{2}\int_{k_2}\frac{1}{k_2^2+m^2}\bigg\}_{{
\cal O}(m^2)}.
\end{equation}
对 $m>0$，我们求得
\begin{align}
d_{\mathbb{I}}^{(1)} = {} & -
\frac{m^2x^2}{128\pi^2}\left[\gamma_{\;\mathrm{E}}-1+\log\left(\frac{x^2}{
8\tau } \right)\right]\left[
1+ \log\left(\frac{\mu^2}{m^2}\right)\right].
\end{align}

这里我们看到，第二项——它是新的重正化尺度 $\mu$ 与裸质量 $m$ 的函数——不过是原理论在
$\overline{MS}$ 方案下单圈的质量重正化贡献 \cite{Kleinert:2001ax}。此外，含流时间的因子与
领头阶贡献 $d_{\mathbb{I}}^{(0)}$ 中 ${\cal O}(m^2)$ 项、即式~\eqref{eq:d1lead}，完全相同。

因此，我们可以简单地把领头阶与次领头阶的项合并为
\begin{align}
d_{\mathbb{I}} = {} & 1 -
\frac{x^2}{8\tau}  +\frac{m_{\mathrm{R}}^2x^2}{4}\left[\gamma_{\;\mathrm{E}} -1
+
\log\left(\frac{x^2}{8\tau}\right)\right]  +{\cal O}(\lambda^2),
\end{align}
其中 $m_{\mathrm{R}}$ 是 $\overline{MS}$ 方案中的重正化质量，由
$m^2_{\mathrm{R}} = Z_m^{-1}m^2$ 给出，而 $Z_m$ 是质量重正化 \cite{Kleinert:2001ax}：
\begin{equation}
Z_m = 1+\frac{\lambda}{16\pi^2}\left[
1+ \log\left(\frac{\mu^2}{m^2}\right)\right]+{\cal O}(\lambda^2).
\end{equation}

这是一个清晰的次领头阶例子，说明非零流时间处理论的发散如何被零流时间处原理论的重正化参数吸收。
换言之，零流时间的重正化理论在非零流时间保持紫外有限。

\subsection{连通贡献}

我们在图~\ref{fig:sopeconn} 与图~\ref{fig:sopehigher} 中分别给出式~\eqref{eq:sope2pt}
中连通威尔逊系数 $d_{\rho^2}(\mu x,\mu^2\tau,mx)$ 的领头阶与次领头阶贡献。
\begin{figure}[htbp]
\includegraphics[width=0.42\textwidth,keepaspectratio=true]{images/dphi1loop}
\caption{\label{fig:sopeconn}
表示威尔逊系数 $d_{\rho^2}$ 领头阶（单圈）贡献的费曼图。费曼图的细节见图~\ref{fig:sope2pt}
的说明。}
\end{figure}
\begin{figure}[htbp]
\begin{minipage}{0.6\textwidth}
\includegraphics[width=0.95\textwidth,keepaspectratio=true]{images/dphi2loopa}
\end{minipage}%
\caption{\label{fig:sopehigher}
表示威尔逊系数 $d_{\rho^2}$ 次领头阶（双圈）贡献的费曼图。关于费曼图的细节，
见图~\ref{fig:sope2pt} 的说明。}
\end{figure}
本节中我们略去那些平凡地并入外部场波函数重正化的图，例如图~\ref{fig:zphi}
所示的单圈例子。只要零流时间的原理论已被重正化，反抵消项就会把这些贡献完全抵消掉。
\begin{figure}[htbp]
\includegraphics[width=0.62\textwidth,keepaspectratio=true]{images/dphizphi}
\caption{\label{fig:zphi}
自然地并入重正化外部传播子的示例图。我们在威尔逊系数的显式计算中不包括这些贡献。关于费曼图的细节，
见图~\ref{fig:sope2pt} 的说明。}
\end{figure}

\subsubsection{领头阶}

我们通过考虑 sOPE 中各算符与两个外部场耦合的矩阵元来提取领头阶的连通威尔逊系数；这会消去一切
非连通贡献。然后，在 ${\cal O}(\lambda)$ 量级上比对各项，即可读出威尔逊系数的单圈贡献
$d_{\rho^2}^{(1)}$。该贡献由如下表达式的小时空行为给出：
\begin{align}
-\lambda d_{\rho^2}^{(1)} {} & \stackrel{x\sim 0}{=}
\Big\{ \big\langle \Omega | {\cal T}\{
\phi(x){} \phi(0)\}
\widetilde{\phi}(p_1) \widetilde{\phi}(p_2) | \Omega \big\rangle^{(1)}
\nonumber\\
{} & - \big\langle \Omega | \rho^2(\tau,0)
\widetilde{\phi}(p_1) \widetilde{\phi}(p_2) | \Omega
\big\rangle^{(1)}\Big\}_{{\cal O}(m^0)}.
\end{align}

相应的费曼积分为
\begin{align}\label{eq:curlyb}
\frac{1}{(p^2_1+m^2)(p_2^2+m^2)}\left\{\frac{1}{2}\int_k
\frac{e^{ik\cdot x} - e^{-(k^2+q^2) \tau }}{(k^2+m^2)(q^2+m^2)}\right\},
\end{align}
其中 $q = k-p_1-p_2$。我们考察花括号内积分的小时空行为并将其展开至 ${\cal O}(m^0)$，
从中提取涂抹威尔逊系数：
\begin{align}\label{eq:drho1}
d_{\rho^2}^{(1)} \stackrel{x\sim 0}{=} \left\{
\frac{1}{2}\int_k
\frac{e^{ik\cdot x} - e^{-(k^2+q^2)\tau }}{(k^2+m^2)(q^2+m^2)}\right\}_{{\cal
O}(m^0)}.
\end{align}

涂抹威尔逊系数必须与外部态无关，才能保证 sOPE 是真正的算符展开。在这个具体情形中，
我们要求 $d_{\rho^2}^{(1)}$ 与外部动量
$p_1$、$p_2$ 以及质量无关。通过对其中一个外部动量求导，
\begin{equation}\label{eq:drhoderiv}
\frac{\mathrm{d}}{\mathrm{d}p_i}d_{\rho^2}^{(1)} =
\int_k \frac{q_i\big[e^{ik\cdot x} -
e^{-(k^2+q^2)\tau}(1+(q^2+m^2)\tau)\big]}{(k^2+m^2)(q^2+m^2)^2},
\end{equation}
我们得到一个收敛的积分。该积分的 $x\rightarrow 0$ 极限现在是良定义的，且仅当流时间与时间空间隔
相关联时才为零。若对外部质量 $m$ 作类似求导，也会得到类似的结果。于是，由量纲分析可知，
选取 $\tau \propto x^2$ 即可保证涂抹威尔逊系数与外部态无关。

然而从物理上讲，我们要求涂抹半径小于非定域算符的时空延展：$s_{\mathrm{rms}}< x$。这一选择保证
sOPE 仍是关于局域算符的展开。直观地说，如果梯度流探测的长度尺度达到非定域算符的量级——如图~\ref%
{fig:sopeexpansion} 所示——那么涂抹算符将不再是（近似）局域的；换言之，sOPE 对原算符而言就是糟糕的
展开。这正是我们在领头阶非连通贡献（式~\eqref{eq:d1lead}）中看到的第三个特征的物理起源：小时间空间隔
要求小流时间以保证收敛。OPE 要求小时间空间隔以获得良好的收敛，因此 sOPE 同样要求小流时间。
\begin{figure}[htbp]
\includegraphics[width=0.62\textwidth,keepaspectratio=true]{images/dphiexpansion}
\caption{\label{fig:sopeexpansion} 在小流时间下（左图），涂抹算符相对于非定域算符的时间空间隔而言是
局域化的。而在大流时间值下（右图），涂抹半径探测到了非定域算符的尺度，此时 sOPE 是对原非定域算符的
糟糕近似。关于费曼图的细节，见图~\ref{fig:sope2pt} 的说明。
}
\end{figure}

有一个互补的视角可以更定量地阐明这一小流时间要求：小流时间展开
\cite{Makino:2014taa,Suzuki:2013gza,Luscher:2011bx}。在此视角下我们看到，
式~\eqref{eq:drhoderiv} 的导数在 $x\rightarrow 0$ 极限下与外部动量无关，
只差流时间的线性项：
\begin{equation}
\frac{\mathrm{d}}{\mathrm{d}p_i}d_{\rho^2}^{(1)} =
\int_k\frac{q_i\big(e^{ik\cdot x} -
1\big)}{(k^2+m^2)(q^2+m^2)^2}+{\cal O}(\tau).
\end{equation}
由此显而易见：只要流时间相对时间空间隔很小，威尔逊系数就与外部态无关（只差流时间的线性项）；
也就是说，sOPE 是在可定量刻画的意义上关于近似局域算符的展开。此外，在一圈以及
${\cal O}(\tau)$ 精度上，流时间依赖会在威尔逊系数与其相应矩阵元的乘积中消去。

由于涂抹威尔逊系数在我们工作的精度上与外部动量和质量无关 \cite{Collins:1984rdg}，
我们可以在式~\eqref{eq:drho1} 中自由地取 $p_1 = p_2 = 0$。对质量展开，得到
\begin{align}
d_{\rho^2}^{(1)} = {} &
\frac{1}{2}\int_k
\frac{e^{ik\cdot x} - e^{-2k^2 \tau }}{(k^2)^2} \nonumber \\
= {} & -\frac{1}{32\pi^2}\left[\gamma_{\;\mathrm{E}} -1 +
\log\left(\frac{x^2}{8\tau}\right)\right].
\end{align}
把它与领头阶贡献（恰为 1）合并，我们有
\begin{equation}\label{eq:drhores}
d_{\rho^2}= 1 +\frac{\lambda}{32\pi^2}\left[\gamma_{\;\mathrm{E}} -1 +
\log\left(\frac{x^2}{8\tau}\right)\right]+{\cal O}(\lambda^2).
\end{equation}

相比之下，$\ms$ 方案中 OPE 的威尔逊系数（记作
$\overline{c}$）为
\begin{equation}\label{eq:cphires}
\overline{c}_{\phi^2} =  1+\frac{\lambda}{32\pi^2} \left[1
+2\gamma_{\;\mathrm{E}}+
\log\left(\frac{\mu^2x^2}{16}\right)\right]+{\cal O}(\lambda^2).
\end{equation}
我们注意到，在领头阶非连通贡献中观察到的三个特征再次出现：涂抹与非涂抹系数具有相同的时空依赖；
流时间作为领头阶调节子出现；以及需要为小时间空间隔选择小流时间，以避免大的对数贡献。
从小流时间展开的视角出发，我们可以确认流时间依赖最终会在所需精度上消去。例如，
$\rho^2(\tau,0)$ 与两个外场耦合的矩阵元为
\begin{align}
\!\!\!\!\big\langle \Omega
|{} & \rho^2(\tau,0)\widetilde{\phi}(p_1) \widetilde{\phi}(p_2)
|\Omega \big\rangle = 1  \nonumber\\
{} &\quad +\frac{\lambda}{32\pi^2}\left[1+\gamma_{\;\mathrm{E}} +
\log\left(2m^2\tau\right)\right]
+{\cal O}(\tau,\lambda^2)
\end{align}
（对流时间足够小的情形）。这里为了简单起见，我们略去了取在壳的外部场。如预期的那样，
该矩阵元与威尔逊系数的乘积在一圈及 ${\cal O}(\tau)$ 精度上与流时间无关：
\begin{align}
d_{\rho^2}\big\langle  \Omega
|{} &\rho^2(\tau,0)\widetilde{\phi}(p_1)\widetilde{\phi}(p_2)
| \Omega \big\rangle =  1 \nonumber\\
{} &
+\frac{\lambda}{32\pi^2}\left[2\gamma_{\;\mathrm{E}} +
\log\left(\frac{m^2x^2}{4}\right)\right]+{\cal O}(\tau,\lambda^2).
\end{align}


\subsubsection{次领头阶}

为确定连通威尔逊系数的次领头阶贡献，必须计入图~\ref{fig:sopehigher}
中的费曼图。图~\ref{fig:2loopexc} 中的双圈图对威尔逊系数 $d_{\rho^2}$
没有贡献，因为它们出现在 ${\cal O}(m^2)$ 量级。
\begin{figure}[htbp]
\begin{minipage}{0.5\textwidth}
\includegraphics[width=0.95\textwidth,keepaspectratio=true]{images/dphi2loopm2}
\end{minipage}
\caption{\label{fig:2loopexc}
出现在 ${\cal O}(m^2)$ 量级、因而对 $d_{\rho^2}$ 无贡献的双圈图。关于费曼图的细节，
见图~\ref{fig:sope2pt} 的说明。}
\end{figure}

观察图~\ref{fig:sopehigher} 中的费曼图，并结合我们对 $d_{\mathbb{I}}$
次领头阶贡献的经验，立即可以看出这些图正是单圈威尔逊系数 $d_{\rho^2}^{(1)}$ 与次领头阶重正化四点顶点的
乘积。当然，这一顶点贡献不是别的，正是重正化耦合常数的次领头阶贡献。于是我们可以相当简单地写出
威尔逊系数
\begin{equation}\label{eq:drho2r}
d_{\rho^2} = 1 +\frac{\lambda_{\mathrm{R}}}{32\pi^2}\left[\gamma_{\;\mathrm{E}}
-1 +
\log\left(\frac{x^2}{8\tau}\right)\right] +{\cal O}(\lambda^3),
\end{equation}
其中 $\lambda_{\mathrm{R}}$ 是重正化耦合。在 $\ms$ 方案中，
重正化耦合由
\begin{equation}\label{eq:lambdar}
\lambda_{\mathrm{R}} = \lambda -
\frac{3\lambda^2}{2(16\pi^2)}\log\left(\frac{\mu^2}{m^2}\right)+{\cal
O}(\lambda^3).
\end{equation}
给出。该量的五圈计算见文献 \cite{Kleinert:2001ax}。

当我们超越领头阶的威尔逊系数（即 $d_{\mathbb{I}}^{(0)}$ 与
$d_{\rho^2}^{(1)}$）时，计算中出现发散的辐射修正，因为所有场相互作用顶点都位于零流时间。
这些发散可以被原理论的重正化参数移除，而且涂抹算符的重正化尺度依赖完全包含在原理论的重正化参数之中。
这一点在预料之中：它源于这样一个事实——零流时间的重正化矩阵元在非零流时间保持有限，且无须进一步
重正化 \cite{Makino:2014sta,Luscher:2011bx}。

最终，对于格点 QCD 中的现实计算，涂抹威尔逊系数的微扰计算必须与强子能标处矩阵元的非微扰确定相结合。
标量 $\phi^4$ 理论在四维并非渐近自由，但通过研究标量场这一简单例子的重正化群方程，
我们可以更细致地理解 sOPE 的数学特征。
